% ============================================================
% UNIDAD 3: LICENCIAS, REPRESENTANTE TÉCNICO Y OFICINA TÉCNICA
% ============================================================
\clearpage
\section*{UNIDAD 3: LICENCIAS DEL PERSONAL Y OFICINA TÉCNICA}
\subsection*{El mecánico, el representante técnico y su ámbito legal}

\begin{recuadro}
\textbf{Objetivo de la unidad:} conocer las licencias del personal de mantenimiento, en especial
la del técnico aviónico; sus requisitos, atribuciones, limitaciones y obligaciones; el rol del
\textbf{representante técnico} de un taller y las funciones de la \textbf{oficina técnica}.
\end{recuadro}

% ------------------------------------------------------------
\section*{1. LAS LICENCIAS DEL PERSONAL AERONÁUTICO}

El Código Aeronáutico establece que el personal aeronáutico necesita una \textbf{licencia o
certificado de idoneidad} expedido por la autoridad para ejercer sus funciones. Los requisitos
siguen el \textbf{Anexo 1} de la OACI. Para el personal que no es tripulante de vuelo, la norma
es la \textbf{RAAC Parte 65}, ``Personal aeronáutico, excepto miembros de la tripulación de
vuelo''.

\begin{center}
\small
\begin{tabular}{|p{0.15\linewidth}|p{0.75\linewidth}|}
\hline
\textbf{Subparte} & \textbf{Licencia o certificado} \\ \hline
D & \textbf{Licencia de Mecánico de Mantenimiento de Aeronave (MMA)}. \\ \hline
J & \textbf{Licencia de Mecánico de Equipos Radioeléctricos de Aeronaves (MERA)}: \textbf{la
licencia del técnico aviónico}. \\ \hline
K & Certificado de competencia de certificador aeronáutico. \\ \hline
L & Certificado de competencia de tareas especiales de mantenimiento. \\ \hline
Otras & Controlador de tránsito aéreo, despachante de aeronave, operador de estación
aeronáutica, jefe de aeródromo, entre otras. \\ \hline
\end{tabular}
\end{center}

% ------------------------------------------------------------
\section*{2. LA LICENCIA MERA}

\subsubsection*{2.1. Requisitos para obtenerla por primera vez}
Con la Resolución ANAC N° 443/2025 se actualizaron los requisitos de las licencias MMA y MERA.
Para la licencia MERA se pide, en síntesis:
\begin{itemize}
  \item Tener \textbf{18 años}.
  \item Ser capaz de hablar, leer, escribir y entender el idioma español.
  \item Acreditar estudios (certificado de nivel primario o secundario).
  \item Poseer el \textbf{título} de Mecánico de Equipos Radioeléctricos de Aeronave o de
        Mecánico de Aviónica otorgado por un Centro de Instrucción certificado bajo la RAAC 147,
        \textbf{o el título de Técnico en Aviónica (o equivalente) emitido por una escuela
        técnica}.
  \item Un \textbf{certificado médico aeronáutico} vigente, de la clase que exija la regulación
        (RAAC Parte 67).
  \item La inscripción vigente en el \textbf{Consejo Profesional de la Ingeniería Aeronáutica y
        Espacial (CPIAyE)}.
  \item DNI, formulario de solicitud, foto y el pago del arancel.
\end{itemize}
El trámite se inicia por la plataforma de \textbf{Trámites a Distancia (TAD)} o en forma
presencial con turno. La autoridad puede verificar el cumplimiento de los requisitos y tomar los
exámenes que correspondan.

\subsubsection*{2.2. Habilitaciones y recategorización}
La licencia se completa con \textbf{habilitaciones}, que indican qué trabajos puede certificar su
titular (en la MMA, categorías A, B o C; en la MERA, la habilitación en aviónica). Para obtenerlas
hay que \textbf{acreditar experiencia de trabajo} con un formulario específico, además de los
requisitos generales.

\subsubsection*{2.3. Atribuciones y limitaciones}
\begin{itemize}
  \item El titular de una licencia \textbf{sin habilitación} puede realizar tareas de
        mantenimiento, pero \textbf{bajo la supervisión} de un titular habilitado con atribuciones
        para esas tareas.
  \item Con la habilitación correspondiente puede realizar y \textbf{certificar} los trabajos
        incluidos en ella, y firmar los registros de mantenimiento dentro de sus atribuciones.
  \item Nunca puede certificar trabajos fuera de las atribuciones de su licencia y habilitación,
        ni sobre tipos de equipos o aeronaves para los que no tiene la instrucción y la experiencia
        necesarias.
\end{itemize}

\subsubsection*{2.4. Validez y obligaciones del titular}
\begin{itemize}
  \item Mantener vigente el \textbf{certificado médico} y no ejercer cuando su aptitud física o
        mental esté disminuida (enfermedad, medicamentos, fatiga).
  \item \textbf{No trabajar bajo los efectos del alcohol o de drogas.}
  \item Llevar consigo la licencia y exhibirla cuando la autoridad lo requiera.
  \item Trabajar con la documentación técnica vigente y las herramientas adecuadas y calibradas.
  \item \textbf{Registrar con veracidad} todo trabajo realizado. Falsear un registro es una falta
        gravísima.
  \item Mantenerse actualizado (cursos, nuevos equipos, cambios normativos).
\end{itemize}

\begin{recuadro}
\textbf{Responsabilidad legal.} Cuando un técnico firma un retorno al servicio, asume la
responsabilidad de que el trabajo se hizo bien y de que la aeronave o el componente quedó
aeronavegable. Si no lo hizo correctamente, puede sufrir \textbf{sanciones administrativas}
(apercibimiento, suspensión o inhabilitación de la licencia por parte de la ANAC) y, si hay daños,
\textbf{responsabilidad civil} (indemnizar) y \textbf{penal} (por ejemplo, por estrago u homicidio
culposo). La firma no es un trámite: es una garantía personal.
\end{recuadro}

% ------------------------------------------------------------
\section*{3. EL REPRESENTANTE TÉCNICO}

\subsubsection*{3.1. Definición}
La RAAC 145 define al \textbf{representante técnico (RT)} como la persona designada por el Taller
Aeronáutico de Reparación (TAR) que:
\begin{itemize}
  \item es \textbf{responsable del retorno al servicio} de cualquier artículo procesado por el
        taller;
  \item asume la \textbf{responsabilidad técnica sobre todas las operaciones} del taller que se
        realizan según la RAAC 145, incluido asegurar que el personal cumpla las regulaciones; y
  \item \textbf{representa al taller ante la autoridad aeronáutica} (ANAC).
\end{itemize}
Junto a él existe el \textbf{gerente responsable}, que tiene la autoridad corporativa (y los
recursos económicos) para asegurar que todo el mantenimiento pueda hacerse según las
regulaciones. Para certificar un taller, la ANAC pide los antecedentes del representante técnico
propuesto y una nota en la que acepta el cargo.

\subsubsection*{3.2. Funciones, derechos y obligaciones}
\begin{center}
\small
\begin{tabular}{|p{0.44\linewidth}|p{0.46\linewidth}|}
\hline
\textbf{Funciones y obligaciones} & \textbf{Derechos y facultades} \\ \hline
\begin{itemize}[leftmargin=*,nosep]
  \item Firmar los retornos al servicio (o designar a quienes lo hagan según el manual).
  \item Mantener actualizado el manual del taller y la lista de capacidades.
  \item Controlar que el personal tenga licencias, habilitaciones e instrucción.
  \item Asegurar herramientas calibradas, documentación vigente e instalaciones adecuadas.
  \item Comunicar a la ANAC cambios, fallas y defectos graves.
  \item Atender las inspecciones y auditorías de la autoridad.
\end{itemize}
&
\begin{itemize}[leftmargin=*,nosep]
  \item Detener un trabajo o negarse a firmar un retorno al servicio si no está conforme, aunque
        lo presione el cliente o la empresa.
  \item Exigir a la empresa los recursos necesarios para trabajar según la norma.
  \item Rechazar partes sin documentación válida.
  \item Ser consultado ante cualquier cambio que afecte la operación técnica del taller.
\end{itemize}
\\ \hline
\end{tabular}
\end{center}

\subsubsection*{3.3. Organigrama funcional de un taller}
\begin{center}
\begin{tikzpicture}[every node/.style={font=\footnotesize, align=center},
  caja/.style={draw=ppAzul, very thick, rounded corners, fill=ppAzul!8, minimum width=3.1cm,
  minimum height=0.9cm}]
  \node[caja] (g) at (0,0) {\textbf{Gerente responsable}\\ recursos y autoridad corporativa};
  \node[caja] (rt) at (-3.3,-1.6) {\textbf{Representante técnico}\\ responsabilidad técnica};
  \node[caja] (cal) at (3.3,-1.6) {\textbf{Calidad}\\ auditoría interna};
  \node[caja] (ot) at (-5.4,-3.3) {\textbf{Oficina técnica}\\ documentación, ingeniería};
  \node[caja] (pl) at (-1.8,-3.3) {\textbf{Planificación}\\ y producción};
  \node[caja] (pa) at (1.8,-3.3) {\textbf{Pañol}\\ y abastecimiento};
  \node[caja] (in) at (5.4,-3.3) {\textbf{Inspección}\\ inspectores};
  \node[caja] (sec) at (-1.8,-4.9) {\textbf{Secciones del taller}\\ aviónica, estructuras, motores};
  \draw[thick] (g) -- (rt); \draw[thick] (g) -- (cal);
  \draw[thick] (rt) -- (ot); \draw[thick] (rt) -- (pl); \draw[thick] (rt) -- (pa);
  \draw[thick] (rt) -- (in); \draw[thick] (pl) -- (sec);
  \draw[thick, dashed] (cal) -- (in);
\end{tikzpicture}
\end{center}
La función de \textbf{calidad} debe ser \textbf{independiente} de la producción: quien inspecciona
no debe depender de quien tiene que terminar el trabajo a tiempo.

% ------------------------------------------------------------
\section*{4. LA OFICINA TÉCNICA}

La \textbf{oficina técnica} es el ``cerebro documental'' de un taller o de un operador. Es uno de
los ámbitos de trabajo más importantes para un técnico con experiencia.

\subsubsection*{4.1. Funciones}
\begin{itemize}
  \item \textbf{Biblioteca técnica:} recibir, controlar y distribuir los manuales en su última
        revisión; retirar las revisiones viejas.
  \item \textbf{Ingeniería:} evaluar las directivas de aeronavegabilidad y los boletines de
        servicio que aplican a cada aeronave; estudiar modificaciones; obtener los
        \textbf{datos aprobados} para reparaciones y alteraciones mayores.
  \item \textbf{Planificación:} armar las órdenes de trabajo y las tarjetas de tareas a partir del
        programa de mantenimiento.
  \item \textbf{Registros:} mantener los historiales de aeronaves, motores y componentes, y el
        control de vencimientos y vidas límite.
  \item \textbf{Relación con la autoridad:} preparar las presentaciones, formularios e informes, y
        atender las inspecciones.
\end{itemize}

\subsubsection*{4.2. Planos, memorias y documentación técnica}
\begin{itemize}
  \item \textbf{Planos:} de conjunto, de detalle, de instalación y \textbf{diagramas eléctricos}
        (esquemáticos, de cableado y de conexionado). Hay que saber leer las vistas, las cotas,
        las tolerancias, los símbolos y la \textbf{lista de materiales}.
  \item \textbf{Memoria técnica:} documento que describe y justifica un trabajo de ingeniería. Para
        una instalación de aviónica suele incluir: descripción del equipo y de la modificación,
        normas aplicables, ubicación y fijación, diagrama eléctrico, \textbf{análisis de carga
        eléctrica} (que el sistema de generación soporte el nuevo consumo), cálculo de \textbf{peso
        y balanceo}, compatibilidad electromagnética, ensayos funcionales y suplemento al manual
        de vuelo si corresponde.
  \item \textbf{Órdenes de trabajo y tarjetas de tareas}, con los pasos, las herramientas, los
        valores de prueba y las firmas de quien ejecuta y de quien inspecciona.
\end{itemize}

\subsubsection*{4.3. La relación con la autoridad de aeronavegabilidad}
Históricamente, la aeronavegabilidad estaba a cargo de la \textbf{Dirección Nacional de
Aeronavegabilidad (DNA)}; hoy esas funciones están en la \textbf{ANAC}, dentro de su Dirección
Nacional de Seguridad Operacional. La oficina técnica se relaciona con ella para:
\begin{itemize}
  \item presentar los trabajos de reparación y alteración mayor y sus datos;
  \item solicitar aprobaciones, inspecciones de aeronavegabilidad y certificados;
  \item informar fallas, defectos y dificultades en servicio;
  \item responder a las observaciones de las auditorías.
\end{itemize}

% ------------------------------------------------------------
\section*{5. ACTIVIDADES DE LA UNIDAD}
\begin{enumerate}
  \item \textbf{Simulación de emisión de la licencia MERA:} armar la carpeta completa de
        solicitud con todos los requisitos (formulario, documentación, título, certificado médico,
        inscripción profesional, pago). Un grupo actúa como ``autoridad'' y la revisa: ¿está
        completa?, ¿qué falta?
  \item \textbf{Atribuciones:} para cinco situaciones del taller, decidir si un técnico con
        licencia sin habilitación, uno con habilitación o el representante técnico puede hacer y
        firmar el trabajo.
  \item \textbf{Organigrama funcional} de un TAR de aviónica, con las funciones, derechos y
        obligaciones del representante técnico.
  \item \textbf{Memoria técnica:} redactar la memoria de la instalación de un transpondedor con
        ADS-B en un avión liviano, incluyendo diagrama eléctrico y análisis de carga eléctrica.
  \item \textbf{Registro y retorno al servicio:} completar el registro de un trabajo de aviónica
        con todos sus datos y la firma de retorno al servicio.
  \item \textbf{Dilema ético:} el cliente apura la entrega de un avión y falta una prueba
        funcional. ¿Qué hace el representante técnico? ¿Qué dice la norma? ¿Qué riesgos hay?
\end{enumerate}
